\pagestyle{empty}
\tight
\begin{tblr}{width=\textwidth,colspec={|Q[c,m,1.9cm]|X[l,m]|},hlines,vlines,hspan=minimal,row{1}={bg=headblue},rowsep=1pt,colsep=3pt}
\SetCell{c,m}\hfil\logo\hfil & \textbf{POLITEKNIK NEGERI MALANG}\par\textbf{JURUSAN TEKNOLOGI INFORMASI}\par\textbf{PROGRAM STUDI: D4 TEKNIK INFORMATIKA} \\
\SetCell[c=2]{c} \textbf{RENCANA PEMBELAJARAN SEMESTER (RPS)} & \\
\end{tblr}
\begin{longtblr}{width=\textwidth,colspec={|Q[l,m,1.9cm]|X[0.85,l,m]|X[1.45,l,m]|X[0.75,l,m]|X[1.15,l,m]|},hlines,vlines,hspan=minimal,rowsep=1pt,colsep=2pt,row{1,3}={bg=headgray,font=\bfseries}}
MATA KULIAH & KODE & BOBOT (SKS)/jam & SEMESTER & TGL. PENYUSUNAN \\
Desain dan Pemrograman Web & RTI253004 & 3 SKS/6 jam & 3 & 6 Agustus 2025 \\
OTORISASI & Dosen Pengembang RPS & Koordinator RMK & \SetCell[c=2]{l} Ka PRODI &  \\
 & Dimas Wahyu Wibowo, S.T., M.T. &  & \SetCell[c=2]{l}{\vspace{8mm}\par Dr. Ely Setyo Astuti, S.T., M.T.} &  \\
\SetCell[r=13]{l} Capaian Pembelajaran (CP) & \SetCell[c=4]{l,bg=headgray,font=\bfseries} Capaian Pembelajaran Lulusan Program Studi (CPL-Prodi) &  &  &  \\
 & CPL02 & \SetCell[c=3]{l} Mampu menerapkan prinsip rekayasa sistem dan perangkat lunak untuk menghasilkan solusi aplikasi multi-platform sesuai kebutuhan. &  &  \\
 & CPL06 & \SetCell[c=3]{l} Mampu merancang, mengimplementasikan, menguji, melakukan deployment, memelihara, serta menjamin mutu perangkat lunak yang menjawab permasalahan dan memenuhi kebutuhan pengguna. &  &  \\
 & CPL07 & \SetCell[c=3]{l} Mampu menerapkan sistem komputasi cerdas dan pengolahan data untuk mendukung penyelesaian masalah komputasi. &  &  \\
 & \SetCell[c=4]{l,bg=headgray,font=\bfseries} Capaian Pembelajaran mata kuliah (CPL-MK) &  &  &  \\
 & CPMK0209 & \SetCell[c=3]{l} Mahasiswa mampu memilih dan menerapkan teknologi pengembangan berbasis platform untuk membangun aplikasi sesuai kebutuhan. &  &  \\
 & CPMK0603 & \SetCell[c=3]{l} Mahasiswa mampu mengimplementasikan dan melakukan deployment aplikasi web, mobile, atau framework sesuai kebutuhan pengguna. &  &  \\
 & CPMK0704 & \SetCell[c=3]{l} Mahasiswa mampu menerapkan berbagai bahasa dan paradigma pemrograman untuk menyelesaikan masalah komputasi. &  &  \\
 & \SetCell[c=4]{l,bg=headgray,font=\bfseries} Kemampuan akhir tiap tahapan belajar (Sub-CPMK) &  &  &  \\
 & SCPMK0704-02201 & \SetCell[c=3]{l} Mahasiswa mampu menerapkan HTML, CSS, JavaScript dasar, dan PHP untuk membangun halaman web. &  &  \\
 & SCPMK0209-02202 & \SetCell[c=3]{l} Mahasiswa mampu memilih mekanisme client-server, form, session, cookie, dan database sesuai kebutuhan aplikasi web. &  &  \\
 & SCPMK0603-02203 & \SetCell[c=3]{l} Mahasiswa mampu mengimplementasikan aplikasi web dinamis dengan PHP dan MySQL serta melakukan debugging. &  &  \\
 & SCPMK0603-02204 & \SetCell[c=3]{l} Mahasiswa mampu menghasilkan aplikasi web sederhana yang dapat dipresentasikan dan diunggah ke hosting. &  &  \\
\SetCell[r=2]{l} Penilaian & \SetCell[c=4]{l} *Nilai CPMK didapat dari agregasi nilai SUB-CPMK yang dibebankan &  &  &  \\
 & \SetCell[c=4]{l}{\tiny\begin{tblr}{width=\linewidth,colspec={|Q[c,m,0.1\linewidth]|Q[c,m,0.16\linewidth]|X[l,m]|X[l,m]|X[l,m]|X[l,m]|X[l,m]|X[l,m]|Q[c,m,0.08\linewidth]|},hlines,vlines,hspan=minimal,rowsep=1pt,colsep=0.7pt,row{1,2}={bg=headgray,font=\bfseries}}
Kode CPMK & Kode SCPMK & \SetCell[c=6]{c} Bobot per Bentuk Penilaian & & & & & & Total Bobot \\
 & & Tugas & Kuis 1 & UTS & Kuis 2 & PBL & UAS & \\
\SetCell[r=1]{c} CPMK0704 & SCPMK0704-02201 & 9\%: HTML, CSS, JavaScript, PHP & 5\%: HTML/CSS/JS & 10\%: modul web dasar & & & & 24.00\% \\
\SetCell[r=1]{c} CPMK0209 & SCPMK0209-02202 & 9\%: form, upload, session, MySQL & & 10\%: modul web dasar & & & & 19.00\% \\
\SetCell[r=2]{c} CPMK0603 & SCPMK0603-02203 & 8\%: CRUD, autentikasi, validasi & & & 5\%: database dan PHP & & 12.50\%: aplikasi web dinamis & 25.50\% \\
 & SCPMK0603-02204 & 8\%: layout, hosting, deployment & & & & 11\%: proyek web dinamis & 12.50\%: aplikasi web dinamis & 31.50\% \\
\SetCell[c=8]{r}\textbf{Total Per Penilaian} & & & & & & & & \textbf{100\%} \\
\end{tblr}} &  &  &  \\
Diskripsi Singkat Mata Kuliah & \SetCell[c=4]{l} Mata kuliah ini membahas dasar desain dan pemrograman web melalui HTML, CSS, JavaScript dasar, PHP, form processing, session/cookie, MySQL, integrasi database, dan hosting. &  &  &  \\
Bahan Kajian / Materi Pembelajaran & \SetCell[c=4]{l}\begin{enumerate}\item HTML, CSS, JavaScript dasar, dan konsep client-server\item PHP, form processing, upload, cookie, dan session\item MySQL dan pemrograman database dengan PHP\item Hosting dan proyek aplikasi web sederhana\end{enumerate} &  &  &  \\
\SetCell[r=2]{l} Pustaka & \SetCell[c=4]{l} \textbf{Utama:}\begin{enumerate}\item Beaird, J. The Principles of Beautiful Web Design.\item Duckett, J. HTML and CSS: Design and Build Websites.\item Pratama, A. PHP Uncover.\end{enumerate} &  &  &  \\
 & \SetCell[c=4]{l} \textbf{Pendukung:} Materi LMS, dokumentasi resmi perangkat lunak, dan jobsheet praktikum. &  &  &  \\
Media Pembelajaran & \SetCell[c=2]{l} \textbf{Software:} OS, web browser, XAMPP/Laragon, MySQL, Visual Studio Code. &  & \SetCell[c=2]{l} \textbf{Hardware:} Komputer/laptop, LCD/proyektor. &  \\
Nama Dosen Pengampu & \SetCell[c=4]{l}\begin{enumerate}\item Priska Choirina, M.Tr.T.\item Muhammad Unggul Pamenang, S.St., M.T.\item Dimas Wahyu Wibowo, S.T., M.T.\item Farid Angga Pribadi, S.Kom., M.Kom.\end{enumerate} &  &  &  \\
Matakuliah Syarat & \SetCell[c=4]{l} - &  &  &  \\
\end{longtblr}
\begin{longtblr}{width=\textwidth,colspec={|Q[c,m,0.06\textwidth]|X[1.35,l,m]|X[1.08,l,m]|X[1.16,l,m]|X[0.7,l,m]|X[1.24,l,m]|X[1.06,l,m]|X[0.98,l,m]|Q[c,m,0.05\textwidth]|},hlines,vlines,hspan=minimal,rowhead=2,row{1,2}={bg=headgreen,font=\bfseries},rowsep=1pt,colsep=1.5pt}
Mgg.\par Ke & Kemampuan Akhir Yang Direncanakan (Sub-CP-MK) & Bahan kajian (Materi Pembelajaran) & Bentuk dan Metode Pembelajaran & Estimasi Waktu & Pengalaman Belajar Mahasiswa & Kriteria \& Bentuk Penilaian & Indikator Penilaian & Bobot Penilaian (\%) \\
(1) & (2) & (3) & (4) & (5) & (6) & (7) & (8) & (9) \\
1 & SCPMK0704-02201 & Konsep web dan HTML. & Modalitas: Blended Learning. Bentuk: Praktikum. Metode: praktik kode, studi kasus, debugging, dan mini project. & 1 x 2 x 50' tatap muka; 1 x 2 x 50' tugas/praktik mandiri. & Mahasiswa mengerjakan aktivitas terarah untuk konsep web dan html dan menunjukkan evidence hasil belajar. & Bentuk penilaian: Konsep web dan HTML. Mengacu rubrik RTM. & Ketepatan konsep; kualitas artefak/praktik; validasi dan dokumentasi. & 1 \\
2 & SCPMK0704-02201 & CSS dan layout. & Modalitas: Blended Learning. Bentuk: Praktikum. Metode: praktik kode, studi kasus, debugging, dan mini project. & 1 x 2 x 50' tatap muka; 1 x 2 x 50' tugas/praktik mandiri. & Mahasiswa mengerjakan aktivitas terarah untuk css dan layout dan menunjukkan evidence hasil belajar. & Bentuk penilaian: CSS dan layout. Mengacu rubrik RTM. & Ketepatan konsep; kualitas artefak/praktik; validasi dan dokumentasi. & 2 \\
3 & SCPMK0704-02201 & JavaScript dasar. & Modalitas: Blended Learning. Bentuk: Praktikum. Metode: praktik kode, studi kasus, debugging, dan mini project. & 1 x 2 x 50' tatap muka; 1 x 2 x 50' tugas/praktik mandiri. & Mahasiswa mengerjakan aktivitas terarah untuk javascript dasar dan menunjukkan evidence hasil belajar. & Bentuk penilaian: JavaScript dasar. Mengacu rubrik RTM. & Ketepatan konsep; kualitas artefak/praktik; validasi dan dokumentasi. & 2 \\
4 & SCPMK0704-02201 & Kuis 1 HTML/CSS/JS. & Modalitas: Blended Learning. Bentuk: Praktikum. Metode: praktik kode, studi kasus, debugging, dan mini project. & 1 x 2 x 50' tatap muka; 1 x 2 x 50' tugas/praktik mandiri. & Mahasiswa mengerjakan aktivitas terarah untuk kuis 1 html/css/js dan menunjukkan evidence hasil belajar. & Bentuk penilaian: Kuis 1 HTML/CSS/JS. Mengacu rubrik RTM. & Ketepatan konsep; kualitas artefak/praktik; validasi dan dokumentasi. & 5 \\
5 & SCPMK0704-02201 & PHP dasar. & Modalitas: Blended Learning. Bentuk: Praktikum. Metode: praktik kode, studi kasus, debugging, dan mini project. & 1 x 2 x 50' tatap muka; 1 x 2 x 50' tugas/praktik mandiri. & Mahasiswa mengerjakan aktivitas terarah untuk php dasar dan menunjukkan evidence hasil belajar. & Bentuk penilaian: PHP dasar. Mengacu rubrik RTM. & Ketepatan konsep; kualitas artefak/praktik; validasi dan dokumentasi. & 2 \\
6 & SCPMK0704-02201 & Percabangan, perulangan, function PHP. & Modalitas: Blended Learning. Bentuk: Praktikum. Metode: praktik kode, studi kasus, debugging, dan mini project. & 1 x 2 x 50' tatap muka; 1 x 2 x 50' tugas/praktik mandiri. & Mahasiswa mengerjakan aktivitas terarah untuk percabangan, perulangan, function php dan menunjukkan evidence hasil belajar. & Bentuk penilaian: Percabangan, perulangan, function PHP. Mengacu rubrik RTM. & Ketepatan konsep; kualitas artefak/praktik; validasi dan dokumentasi. & 2 \\
7 & SCPMK0209-02202 & Form processing dan upload. & Modalitas: Blended Learning. Bentuk: Praktikum. Metode: praktik kode, studi kasus, debugging, dan mini project. & 1 x 2 x 50' tatap muka; 1 x 2 x 50' tugas/praktik mandiri. & Mahasiswa mengerjakan aktivitas terarah untuk form processing dan upload dan menunjukkan evidence hasil belajar. & Bentuk penilaian: Form processing dan upload. Mengacu rubrik RTM. & Ketepatan konsep; kualitas artefak/praktik; validasi dan dokumentasi. & 3 \\
8 & SCPMK0704-02201, SCPMK0209-02202 & UTS modul web dasar. & Modalitas: Blended Learning. Bentuk: Praktikum. Metode: praktik kode, studi kasus, debugging, dan mini project. & 1 x 2 x 50' tatap muka; 1 x 2 x 50' tugas/praktik mandiri. & Mahasiswa mengerjakan aktivitas terarah untuk uts modul web dasar dan menunjukkan evidence hasil belajar. & Bentuk penilaian: UTS modul web dasar. Mengacu rubrik RTM. & Ketepatan konsep; kualitas artefak/praktik; validasi dan dokumentasi. & 20 \\
9 & SCPMK0209-02202 & Cookie dan session. & Modalitas: Blended Learning. Bentuk: Praktikum. Metode: praktik kode, studi kasus, debugging, dan mini project. & 1 x 2 x 50' tatap muka; 1 x 2 x 50' tugas/praktik mandiri. & Mahasiswa mengerjakan aktivitas terarah untuk cookie dan session dan menunjukkan evidence hasil belajar. & Bentuk penilaian: Cookie dan session. Mengacu rubrik RTM. & Ketepatan konsep; kualitas artefak/praktik; validasi dan dokumentasi. & 3 \\
10 & SCPMK0209-02202 & MySQL dan query dasar. & Modalitas: Blended Learning. Bentuk: Praktikum. Metode: praktik kode, studi kasus, debugging, dan mini project. & 1 x 2 x 50' tatap muka; 1 x 2 x 50' tugas/praktik mandiri. & Mahasiswa mengerjakan aktivitas terarah untuk mysql dan query dasar dan menunjukkan evidence hasil belajar. & Bentuk penilaian: MySQL dan query dasar. Mengacu rubrik RTM. & Ketepatan konsep; kualitas artefak/praktik; validasi dan dokumentasi. & 3 \\
11 & SCPMK0603-02203 & Koneksi PHP-MySQL dan CRUD. & Modalitas: Blended Learning. Bentuk: Praktikum. Metode: praktik kode, studi kasus, debugging, dan mini project. & 1 x 2 x 50' tatap muka; 1 x 2 x 50' tugas/praktik mandiri. & Mahasiswa mengerjakan aktivitas terarah untuk koneksi php-mysql dan crud dan menunjukkan evidence hasil belajar. & Bentuk penilaian: Koneksi PHP-MySQL dan CRUD. Mengacu rubrik RTM. & Ketepatan konsep; kualitas artefak/praktik; validasi dan dokumentasi. & 4 \\
12 & SCPMK0603-02203 & Kuis 2 database dan PHP. & Modalitas: Blended Learning. Bentuk: Praktikum. Metode: praktik kode, studi kasus, debugging, dan mini project. & 1 x 2 x 50' tatap muka; 1 x 2 x 50' tugas/praktik mandiri. & Mahasiswa mengerjakan aktivitas terarah untuk kuis 2 database dan php dan menunjukkan evidence hasil belajar. & Bentuk penilaian: Kuis 2 database dan PHP. Mengacu rubrik RTM. & Ketepatan konsep; kualitas artefak/praktik; validasi dan dokumentasi. & 5 \\
13 & SCPMK0603-02203 & Autentikasi dasar dan validasi. & Modalitas: Blended Learning. Bentuk: Praktikum. Metode: praktik kode, studi kasus, debugging, dan mini project. & 1 x 2 x 50' tatap muka; 1 x 2 x 50' tugas/praktik mandiri. & Mahasiswa mengerjakan aktivitas terarah untuk autentikasi dasar dan validasi dan menunjukkan evidence hasil belajar. & Bentuk penilaian: Autentikasi dasar dan validasi. Mengacu rubrik RTM. & Ketepatan konsep; kualitas artefak/praktik; validasi dan dokumentasi. & 4 \\
14 & SCPMK0603-02204 & Template dan layout responsif. & Modalitas: Blended Learning. Bentuk: Praktikum. Metode: praktik kode, studi kasus, debugging, dan mini project. & 1 x 2 x 50' tatap muka; 1 x 2 x 50' tugas/praktik mandiri. & Mahasiswa mengerjakan aktivitas terarah untuk template dan layout responsif dan menunjukkan evidence hasil belajar. & Bentuk penilaian: Template dan layout responsif. Mengacu rubrik RTM. & Ketepatan konsep; kualitas artefak/praktik; validasi dan dokumentasi. & 4 \\
15 & SCPMK0603-02204 & Hosting dan deployment dasar. & Modalitas: Blended Learning. Bentuk: Praktikum. Metode: praktik kode, studi kasus, debugging, dan mini project. & 1 x 2 x 50' tatap muka; 1 x 2 x 50' tugas/praktik mandiri. & Mahasiswa mengerjakan aktivitas terarah untuk hosting dan deployment dasar dan menunjukkan evidence hasil belajar. & Bentuk penilaian: Hosting dan deployment dasar. Mengacu rubrik RTM. & Ketepatan konsep; kualitas artefak/praktik; validasi dan dokumentasi. & 4 \\
16 & SCPMK0603-02204 & PBL proyek web dinamis. & Modalitas: Blended Learning. Bentuk: Praktikum. Metode: praktik kode, studi kasus, debugging, dan mini project. & 1 x 2 x 50' tatap muka; 1 x 2 x 50' tugas/praktik mandiri. & Mahasiswa mengerjakan aktivitas terarah untuk pbl proyek web dinamis dan menunjukkan evidence hasil belajar. & Bentuk penilaian: PBL proyek web dinamis. Mengacu rubrik RTM. & Ketepatan konsep; kualitas artefak/praktik; validasi dan dokumentasi. & 11 \\
17 & SCPMK0603-02203, SCPMK0603-02204 & UAS aplikasi web dinamis. & Modalitas: Blended Learning. Bentuk: Praktikum. Metode: praktik kode, studi kasus, debugging, dan mini project. & 1 x 2 x 50' tatap muka; 1 x 2 x 50' tugas/praktik mandiri. & Mahasiswa mengerjakan aktivitas terarah untuk uas aplikasi web dinamis dan menunjukkan evidence hasil belajar. & Bentuk penilaian: UAS aplikasi web dinamis. Mengacu rubrik RTM. & Ketepatan konsep; kualitas artefak/praktik; validasi dan dokumentasi. & 25 \\
\end{longtblr}
